\documentclass[aspectratio=169]{beamer}
\input{header}
\coursecode{JKSSB}

\title{Excel Formulas}
\subtitle{Lesson 25 --- References, Operators and Copying}
\author{JKSSB Computer Awareness}
\date{\today}

\begin{document}
\frame{\titlepage}

\begin{frame}{What a Formula Does}
  \begin{itemize}
    \item A \key{formula} calculates a new value from other cells, from
      constants, or from functions.
    \item It is typed into the formula bar, and its result appears in the
      cell.
    \item Stored in the cell is the \key{formula}; shown on the sheet is the
      computed result.
    \item Excel recalculates the result automatically when a referenced cell
      changes.
  \end{itemize}
  \begin{block}{The one idea}
    A formula is an instruction that tells Excel how to calculate the value of
    a cell.
  \end{block}
\end{frame}

\begin{frame}{A Formula Must Start with \texttt{=}}
  \begin{itemize}
    \item Every formula begins with the equal sign \key{\texttt{=}}.
    \item \texttt{=A1+A2} is calculated; \texttt{A1+A2} entered without the
      \texttt{=} is stored as ordinary text.
    \item The same symbol also acts as a comparison operator that tests
      equality.
  \end{itemize}
  \begin{alertblock}{Trap alert (TRAP-30, PYQ-38)}
    A common wrong answer says a formula may start with \texttt{+}, \texttt{\%},
    or \texttt{@}. The exam rule is: a \key{formula must start with
    \texttt{=}}.
  \end{alertblock}
\end{frame}

\begin{frame}{Formula Anatomy}
  \begin{itemize}
    \item Take the formula \texttt{=A1*B1+100}.
    \item Leading \texttt{=}: starts the formula.
    \item Operands: cell references (\texttt{A1}, \texttt{B1}) and constants
      (\texttt{100}).
    \item Operators: \texttt{*} (multiply) and \texttt{+} (add).
  \end{itemize}
  \begin{block}{Read it out}
    "Take the value in A1, multiply by B1, then add 100."
  \end{block}
\end{frame}

\begin{frame}{Arithmetic Operators}
  \begin{center}
    \begin{tabular}{p{0.12\textwidth}p{0.62\textwidth}}
      \toprule
      \textbf{Operator} & \textbf{Meaning} \\
      \midrule
      \texttt{+} & Addition \\
      \texttt{-} & Subtraction \\
      \texttt{*} & Multiplication \\
      \texttt{/} & Division \\
      \texttt{\textasciicircum} & Exponent (power) \\
      \texttt{\%} & Percent \\
      \bottomrule
    \end{tabular}
  \end{center}
  \vspace{0.25cm}
  \begin{alertblock}{Exact symbols}
    Use \texttt{*} for multiplication and \texttt{/} for division;
    \texttt{x} and $\div$ are not Excel operators.
    \texttt{=2\textasciicircum3} gives 8.
  \end{alertblock}
\end{frame}

\begin{frame}{Comparison Operators}
  \begin{center}
    \begin{tabular}{p{0.18\textwidth}p{0.56\textwidth}}
      \toprule
      \textbf{Operator} & \textbf{Meaning} \\
      \midrule
      \texttt{=} & Equal to \\
      \texttt{>} & Greater than \\
      \texttt{<} & Less than \\
      \texttt{>=} & Greater than or equal to \\
      \texttt{<=} & Less than or equal to \\
      \texttt{<>} & Not equal to \\
      \bottomrule
    \end{tabular}
  \end{center}
  \vspace{0.25cm}
  \muted{A comparison formula returns TRUE or FALSE, not a number.}
\end{frame}

\begin{frame}{Text and Reference Operators}
  \begin{itemize}
    \item \key{\texttt{\&}} (concatenation) joins text: \texttt{="Total " \& A1}.
    \item \key{\texttt{:}} (range): \texttt{A1:A10} means all cells from A1
      to A10.
    \item \key{\texttt{,}} (union) joins two ranges into one list.
    \item A space (intersection) gives the cells common to two ranges.
  \end{itemize}
  \begin{block}{Why operators matter}
    An operator tells Excel what to do with the operands on either side of it.
  \end{block}
\end{frame}

\begin{frame}{Order of Evaluation (Precedence)}
  \begin{itemize}
    \item Parentheses are evaluated first.
    \item Then percent, then exponent, then multiply/divide, then add/subtract.
    \item Then \texttt{\&} (concatenation).
    \item Then the comparison operators.
    \item Use parentheses to force an order: \texttt{=(A1+A2)*B1}.
  \end{itemize}
  \begin{alertblock}{Trap alert}
    \texttt{=2+3*4} is 14, not 20, because multiplication is done before
    addition.
  \end{alertblock}
\end{frame}

\begin{frame}{Cell References Are Addresses}
  \begin{itemize}
    \item A reference names a cell by its column letter and row number.
    \item \texttt{A1} is column A, row 1; \texttt{B5} is column B, row 5.
    \item A formula uses the \key{current value} of the referenced cell.
    \item Change the referenced cell and the formula result updates at once.
  \end{itemize}
\end{frame}

\begin{frame}{Relative References}
  \begin{itemize}
    \item \key{\texttt{A1}} is a \key{relative} reference.
    \item It is stored as an offset from the formula's own cell.
    \item Copy the formula and the reference shifts by the same offset.
  \end{itemize}
  \begin{exampleblock}{Example}
    \texttt{=A1} placed in B1 and copied to B2 becomes \texttt{=A2}.
  \end{exampleblock}
\end{frame}

\begin{frame}{Absolute References}
  \begin{itemize}
    \item \key{\texttt{\$A\$1}} is an \key{absolute} reference.
    \item The dollar sign locks \key{both the column and the row}.
    \item Copying the formula leaves \texttt{\$A\$1} unchanged.
    \item Use it for a constant such as a fixed rate or a fixed limit.
  \end{itemize}
  \begin{block}{Remember}
    Absolute means anchored: it never moves when the formula is copied.
  \end{block}
\end{frame}

\begin{frame}{Mixed References}
  \begin{itemize}
    \item \key{\texttt{\$A1}}: column locked, row relative.
    \item \key{\texttt{A\$1}}: row locked, column relative.
    \item The \texttt{\$} locks whatever it directly precedes.
    \item \texttt{\$} before a letter locks the column; \texttt{\$} before a
      number locks the row.
  \end{itemize}
\end{frame}

\begin{frame}{Copying Formulas: What Changes}
  \begin{center}
    \begin{tabular}{p{0.20\textwidth}p{0.28\textwidth}p{0.30\textwidth}}
      \toprule
      \textbf{Reference} & \textbf{Locks} & \textbf{B1 copied to C2} \\
      \midrule
      \texttt{A1} & nothing & \texttt{B2} \\
      \texttt{\$A\$1} & column and row & \texttt{\$A\$1} \\
      \texttt{\$A1} & column & \texttt{\$A2} \\
      \texttt{A\$1} & row & \texttt{B\$1} \\
      \bottomrule
    \end{tabular}
  \end{center}
  \vspace{0.25cm}
  \muted{Only the unlocked parts of a reference shift when the formula is
  copied.}
\end{frame}

\begin{frame}{Worked Example: A Fixed Tax Rate}
  \begin{itemize}
    \item Prices are listed in column A; the tax rate is fixed in
      \texttt{\$B\$1}.
    \item Formula in C2: \texttt{=A2*\$B\$1}.
    \item Copy down --- C3 is \texttt{=A3*\$B\$1}, C4 is \texttt{=A4*\$B\$1};
      the rate stays \texttt{\$B\$1}.
    \item If the formula were \texttt{=A2*B1}, copying down would wrongly move
      the rate to B2, B3.
  \end{itemize}
\end{frame}

\begin{frame}{Common Formula Errors}
  \begin{center}
    \begin{tabular}{p{0.18\textwidth}p{0.62\textwidth}}
      \toprule
      \textbf{Error} & \textbf{Meaning} \\
      \midrule
      \texttt{\#DIV/0!} & Divide by zero or by an empty cell \\
      \texttt{\#VALUE!} & Wrong type of operand (text where a number is needed) \\
      \texttt{\#REF!} & A referenced cell was deleted, so the reference is invalid \\
      \texttt{\#NAME?} & Unknown function, or a misspelled name/text entry \\
      \texttt{\#N/A} & The value is not available (for example a failed lookup) \\
      \texttt{\#NULL!} & Two ranges that do not intersect \\
      \bottomrule
    \end{tabular}
  \end{center}
\end{frame}

\begin{frame}{Trap Alert: References and Errors}
  \begin{itemize}
    \item TRAP-30: a formula must start with \key{\texttt{=}} (PYQ-38).
    \item \texttt{\$} before the letter locks the column; \texttt{\$} before
      the number locks the row.
    \item Relative references change on copy; absolute references do not.
    \item \texttt{\#REF!} means a broken reference;
      \texttt{\#DIV/0!} means division by zero.
  \end{itemize}
\end{frame}

\begin{frame}{Lesson 25: Recall}
  \begin{itemize}
    \item A formula begins with \texttt{=} and calculates a cell's value.
    \item Arithmetic operators: \texttt{+ - * / \textasciicircum \%};
      comparison operators include \texttt{<>} for "not equal".
    \item \texttt{\&} joins text; \texttt{:} defines a range;
      \texttt{,} unions ranges.
    \item \texttt{A1} is relative, \texttt{\$A\$1} is absolute,
      \texttt{\$A1} and \texttt{A\$1} are mixed.
    \item On copy, only the unlocked parts of a reference change.
    \item Know the error codes: \texttt{\#DIV/0!}, \texttt{\#VALUE!},
      \texttt{\#REF!}, \texttt{\#NAME?}, \texttt{\#N/A}, \texttt{\#NULL!}.
  \end{itemize}
\end{frame}
\end{document}
